\begin{abstract}
Semi-supervised segmentation typically assumes that an augmented image remains compatible with an unchanged target. This assumption breaks for through-plane MRI perturbations: a finite-support observation combines signals from neighboring slices and therefore changes the tissue composition represented at each pixel. We present \method, a training-only framework that models such perturbations as paired image--target observations. A shared convex slice-profile operator transforms both neighboring MRI slices and their exact or teacher-derived occupancy maps, producing acquisition-aligned fractional supervision for labeled and unlabeled data. Rather than sampling slice profiles heuristically, \method designs a global sampling distribution by maximizing the worst patient--axial-stratum occupancy signal under moment, image-residual, symmetry, entropy, and density constraints. This constrained design preserves the perturbation budget of a predefined parent family while allocating more exposure to informative through-plane transitions. The resulting method is optimized within a standard EMA teacher--student framework and retains the original single-slice 2D network at inference, adding neither parameters nor latency. We evaluate the method on binary prostate and multi-class cardiac MRI segmentation using matched baselines, a compact controlled ablation, patient-level statistics, and cross-dataset transfer without validation- or test-driven profile tuning.
\end{abstract}
